\section{Conclusion}\label{sec:conclusion}

Sparse autoencoders can reconstruct language-model activations well while failing to preserve all of the information
and causal structure that the model uses. In GPT-2 Small, reconstruction quality depended on layer and architecture,
and within a layer it robustly predicted how much splicing an SAE damaged next-token prediction, but it did not robustly
predict truth-information loss, persona information was lost even by SAEs that reconstructed well, and at later sites
the reconstruction retained only about two-thirds of a located causal effect. Improvement is possible but
level-specific. Scaling data and width substantially improved reconstruction and behavioural faithfulness on held-out
data without a confirmed gain in truth information, and behaviour-targeted training improved behaviour at a cost in
reconstruction and causal preservation. In Gemma-2-2B, public Gemma Scope SAEs again produced more behavioural damage
than norm-matched random errors. No single metric establishes that an SAE is faithful. Reconstruction quality is
necessary evidence, but it is insufficient by itself, and future SAE evaluations should report reconstruction,
information, behavioural and causal evidence separately, with the reader, the reference constraint and the evidence
status stated for each.
